\documentclass[11pt]{article}
\usepackage[margin=1in]{geometry}
\usepackage{xcolor}
\usepackage{tcolorbox}
\usepackage{hyperref}
\usepackage{enumitem}
\setlist{noitemsep,topsep=0pt}
\usepackage{booktabs}
\usepackage{array}
\usepackage{amsmath}

% Define OSU orange and AI blue colors
\definecolor{osaorange}{RGB}{220,115,38}
\definecolor{aiblue}{RGB}{30,144,255}

% Configure hyperref colors
\hypersetup{
    colorlinks=true,
    linkcolor=blue,
    urlcolor=blue,
    citecolor=blue
}

\title{}
\author{}
\date{}

\begin{document}

% Orange header box with black outline
\begin{tcolorbox}[colback=osaorange,colframe=black,boxrule=2pt,arc=0pt,left=6pt,right=6pt,top=10pt,bottom=10pt]
\centering
\textcolor{white}{\textbf{\Large In-Class Worksheet}}\\
\textcolor{white}{\textbf{\large Probability and Conditional Probability}}
\end{tcolorbox}

\vspace{8pt}

\noindent\textbf{Name:} \rule{8cm}{0.4pt} \hfill \textbf{Date:} \rule{4cm}{0.4pt}\\[0.5ex]
\textbf{Course:} H524 Introduction to Biostatistics\\
\textbf{Topic:} Probability, Conditional Probability, and Diagnostic Testing

\vspace{8pt}

\section*{Instructions}
Work with a partner to solve the following problems. Use the 2x2 table approach with natural frequencies (concrete counts) rather than formulas whenever possible. Show all your work.

\vspace{12pt}

\section{Problem 1: Mammogram Screening (Based on Gigerenzer's Study)}

A 45-year-old woman with no symptoms or family history of breast cancer has a positive mammogram. What is the probability that she actually has breast cancer?

\vspace{6pt}

\textbf{Given information:}
\begin{itemize}
\item The prevalence of breast cancer in this age group is 1.2\%
\item If a woman has breast cancer, the probability is 85\% that she will have a positive mammogram (sensitivity)
\item If a woman does not have breast cancer, the probability is 9\% that she will still have a positive mammogram (false positive rate)
\end{itemize}

\vspace{12pt}

\textbf{Part A:} Set up a 2x2 table assuming 10,000 women are screened. Fill in all cells:

\vspace{12pt}

\begin{center}
\begin{tabular}{|l|c|c|c|}
\hline
& \textbf{Cancer +} & \textbf{Cancer -} & \textbf{Total} \\
\hline
\textbf{Test +} & & & \\[2ex]
\hline
\textbf{Test -} & & & \\[2ex]
\hline
\textbf{Total} & & & 10,000 \\[2ex]
\hline
\end{tabular}
\end{center}

\vspace{12pt}

\textbf{Work space for calculations:}

\vspace{4cm}

\textbf{Part B:} Calculate the Positive Predictive Value (PPV): $\Pr(\text{Cancer} \mid \text{Test+})$

\vspace{2cm}

\textbf{Part C:} Interpret your result. Why is the PPV so much lower than the sensitivity?

\vspace{2.5cm}

\newpage

\section{Problem 2: Strep Throat Testing}

A rapid strep test has 85\% sensitivity and 95\% specificity. Among patients presenting with sore throat, 15\% actually have strep throat.

\vspace{12pt}

\textbf{Part A:} Create a 2x2 table for 1,000 patients with sore throat:

\vspace{12pt}

\begin{center}
\begin{tabular}{|l|c|c|c|}
\hline
& \textbf{Strep +} & \textbf{Strep -} & \textbf{Total} \\
\hline
\textbf{Test +} & & & \\[2ex]
\hline
\textbf{Test -} & & & \\[2ex]
\hline
\textbf{Total} & & & 1,000 \\[2ex]
\hline
\end{tabular}
\end{center}

\vspace{12pt}

\textbf{Part B:} Calculate the Negative Predictive Value (NPV): $\Pr(\text{No Strep} \mid \text{Test-})$

\vspace{2.5cm}

\textbf{Part C:} What is the probability that a patient with a negative test actually has strep (false negative)?

\vspace{2cm}

\section{Problem 3: Smoking and Heart Disease}

In a study of 1,000 people, there were 200 smokers and 800 non-smokers. During the 10-year follow-up, 30 smokers and 20 non-smokers developed heart disease.

\vspace{12pt}

\textbf{Part A:} Create a 2x2 table:

\vspace{12pt}

\begin{center}
\begin{tabular}{|l|c|c|c|}
\hline
& \textbf{Heart Disease} & \textbf{No Disease} & \textbf{Total} \\
\hline
\textbf{Smokers} & & & \\[2ex]
\hline
\textbf{Non-smokers} & & & \\[2ex]
\hline
\textbf{Total} & & & 1,000 \\[2ex]
\hline
\end{tabular}
\end{center}

\vspace{12pt}

\textbf{Part B:} Calculate the risk of heart disease in smokers:

\vspace{1.5cm}

\textbf{Part C:} Calculate the risk of heart disease in non-smokers:

\vspace{1.5cm}

\textbf{Part D:} Calculate the Relative Risk (RR):

\vspace{1.5cm}

\textbf{Part E:} Interpret the relative risk in plain language:

\vspace{2cm}

\section{Problem 4: Binomial Probability}

A new treatment has a 60\% success rate.

\vspace{12pt}

\textbf{Part A:} If 10 patients receive the treatment, what is the probability that at least 8 of them respond successfully? Use R or a binomial probability table.

\vspace{3cm}

\textbf{Part B:} What is the expected number of successful responses out of 10 patients?

\vspace{1.5cm}

\section{Problem 5: Normal Distribution}

Cholesterol levels in adults are normally distributed with mean 200 mg/dL and standard deviation 40 mg/dL.

\vspace{12pt}

\textbf{Part A:} What proportion of adults have cholesterol levels between 180 and 220 mg/dL?

\vspace{3cm}

\textbf{Part B:} What cholesterol level corresponds to the 90th percentile?

\vspace{3cm}

\end{document}
